%%%%%%%%%%%%%%%%%%%%%%%%%%%%%%%%%%%%%%%%%%%%%%%%%%%%%%%%%%%%
% 14.8 LIMIT LAWS
% The Composition of Advanced Mathematics
%%%%%%%%%%%%%%%%%%%%%%%%%%%%%%%%%%%%%%%%%%%%%%%%%%%%%%%%%%%%

\section{Limit Laws}
\label{sec:limit-laws}

\prerequisite{
Functions, algebraic manipulation, trigonometric functions,
exponential and logarithmic functions, sequences, inequalities,
and basic calculus notation.
}

\learningobjectives{
By the end of this reference section, the reader should be able to:
\begin{itemize}
    \item interpret the notation for finite and infinite limits;
    \item distinguish two-sided and one-sided limits;
    \item use constant, identity, sum, difference, product, quotient,
          power, and root laws;
    \item evaluate polynomial and rational-function limits;
    \item understand when direct substitution is valid;
    \item recognize indeterminate forms;
    \item simplify removable discontinuities algebraically;
    \item use conjugates in radical limits;
    \item use fundamental trigonometric limits;
    \item apply the squeeze theorem;
    \item understand limits of compositions;
    \item recognize continuity-based limit evaluation;
    \item evaluate limits at infinity;
    \item compare polynomial growth rates;
    \item compare logarithmic, polynomial, exponential, and factorial growth;
    \item recognize horizontal asymptotes;
    \item distinguish infinite limits from limits at infinity;
    \item use dominant-term analysis;
    \item interpret sequence limits;
    \item use monotone and bounded sequence principles conceptually;
    \item recognize asymptotic equivalence;
    \item use Big-O and little-o notation conceptually;
    \item recognize common special limits;
    \item understand the connection between limits and continuity;
    \item avoid common errors involving infinity and indeterminate forms.
\end{itemize}
}

%%%%%%%%%%%%%%%%%%%%%%%%%%%%%%%%%%%%%%%%%%%%%%%%%%%%%%%%%%%%
\subsection{Meaning of a Limit}
\label{subsec:meaning-of-limit-reference}
%%%%%%%%%%%%%%%%%%%%%%%%%%%%%%%%%%%%%%%%%%%%%%%%%%%%%%%%%%%%

The notation

\[
\boxed{
\lim_{x\to a}f(x)=L
}
\]

is read:

\[
\boxed{
\text{the limit as }x\text{ approaches }a
\text{ of }f(x)\text{ is }L.
}
\]

Informally, this means that \(f(x)\) can be made arbitrarily close to \(L\)
by taking \(x\) sufficiently close to \(a\), without requiring \(x=a\).

%%%%%%%%%%%%%%%%%%%%%%%%%%%%%%%%%%%%%%%%%%%%%%%%%%%%%%%%%%%%
\subsection{Formal Epsilon--Delta Definition}
\label{subsec:epsilon-delta-limit-reference}
%%%%%%%%%%%%%%%%%%%%%%%%%%%%%%%%%%%%%%%%%%%%%%%%%%%%%%%%%%%%

The statement

\[
\boxed{
\lim_{x\to a}f(x)=L
}
\]

means that for every

\[
\varepsilon>0,
\]

there exists

\[
\delta>0
\]

such that

\[
0<|x-a|<\delta
\]

implies

\[
\boxed{
|f(x)-L|<\varepsilon.
}
\]

The symbols are read:

\[
\varepsilon
=
\text{epsilon},
\]

and

\[
\delta
=
\text{delta}.
\]

%%%%%%%%%%%%%%%%%%%%%%%%%%%%%%%%%%%%%%%%%%%%%%%%%%%%%%%%%%%%
\subsection{Two-Sided Limits}
\label{subsec:two-sided-limit-reference}
%%%%%%%%%%%%%%%%%%%%%%%%%%%%%%%%%%%%%%%%%%%%%%%%%%%%%%%%%%%%

A two-sided limit requires agreement from both directions.

\[
\boxed{
\lim_{x\to a}f(x)=L
}
\]

exists if and only if

\[
\boxed{
\lim_{x\to a^-}f(x)
=
\lim_{x\to a^+}f(x)
=
L.
}
\]

%%%%%%%%%%%%%%%%%%%%%%%%%%%%%%%%%%%%%%%%%%%%%%%%%%%%%%%%%%%%
\subsection{Left-Hand Limits}
\label{subsec:left-hand-limit-reference}
%%%%%%%%%%%%%%%%%%%%%%%%%%%%%%%%%%%%%%%%%%%%%%%%%%%%%%%%%%%%

The notation

\[
\boxed{
\lim_{x\to a^-}f(x)
}
\]

means that \(x\) approaches \(a\) from values less than \(a\).

The superscript minus indicates approach from the left.

%%%%%%%%%%%%%%%%%%%%%%%%%%%%%%%%%%%%%%%%%%%%%%%%%%%%%%%%%%%%
\subsection{Right-Hand Limits}
\label{subsec:right-hand-limit-reference}
%%%%%%%%%%%%%%%%%%%%%%%%%%%%%%%%%%%%%%%%%%%%%%%%%%%%%%%%%%%%

The notation

\[
\boxed{
\lim_{x\to a^+}f(x)
}
\]

means that \(x\) approaches \(a\) from values greater than \(a\).

The superscript plus indicates approach from the right.

%%%%%%%%%%%%%%%%%%%%%%%%%%%%%%%%%%%%%%%%%%%%%%%%%%%%%%%%%%%%
\subsection{When a Two-Sided Limit Does Not Exist}
\label{subsec:two-sided-limit-dne}
%%%%%%%%%%%%%%%%%%%%%%%%%%%%%%%%%%%%%%%%%%%%%%%%%%%%%%%%%%%%

If

\[
\lim_{x\to a^-}f(x)
\neq
\lim_{x\to a^+}f(x),
\]

then

\[
\boxed{
\lim_{x\to a}f(x)
\text{ does not exist}.
}
\]

This commonly occurs at jump discontinuities.

%%%%%%%%%%%%%%%%%%%%%%%%%%%%%%%%%%%%%%%%%%%%%%%%%%%%%%%%%%%%
\subsection{Constant Law}
\label{subsec:constant-limit-law}
%%%%%%%%%%%%%%%%%%%%%%%%%%%%%%%%%%%%%%%%%%%%%%%%%%%%%%%%%%%%

For constant \(c\),

\[
\boxed{
\lim_{x\to a}c
=
c.
}
\]

%%%%%%%%%%%%%%%%%%%%%%%%%%%%%%%%%%%%%%%%%%%%%%%%%%%%%%%%%%%%
\subsection{Identity Law}
\label{subsec:identity-limit-law}
%%%%%%%%%%%%%%%%%%%%%%%%%%%%%%%%%%%%%%%%%%%%%%%%%%%%%%%%%%%%

\[
\boxed{
\lim_{x\to a}x
=
a.
}
\]

More generally,

\[
\boxed{
\lim_{x\to a}x^n
=
a^n
}
\]

for positive integers \(n\).

%%%%%%%%%%%%%%%%%%%%%%%%%%%%%%%%%%%%%%%%%%%%%%%%%%%%%%%%%%%%
\subsection{Constant Multiple Law}
\label{subsec:constant-multiple-limit-law}
%%%%%%%%%%%%%%%%%%%%%%%%%%%%%%%%%%%%%%%%%%%%%%%%%%%%%%%%%%%%

If

\[
\lim_{x\to a}f(x)=L,
\]

then for constant \(c\),

\[
\boxed{
\lim_{x\to a}
cf(x)
=
cL.
}
\]

%%%%%%%%%%%%%%%%%%%%%%%%%%%%%%%%%%%%%%%%%%%%%%%%%%%%%%%%%%%%
\subsection{Sum Law}
\label{subsec:sum-limit-law}
%%%%%%%%%%%%%%%%%%%%%%%%%%%%%%%%%%%%%%%%%%%%%%%%%%%%%%%%%%%%

If

\[
\lim_{x\to a}f(x)=L
\]

and

\[
\lim_{x\to a}g(x)=M,
\]

then

\[
\boxed{
\lim_{x\to a}
[f(x)+g(x)]
=
L+M.
}
\]

%%%%%%%%%%%%%%%%%%%%%%%%%%%%%%%%%%%%%%%%%%%%%%%%%%%%%%%%%%%%
\subsection{Difference Law}
\label{subsec:difference-limit-law}
%%%%%%%%%%%%%%%%%%%%%%%%%%%%%%%%%%%%%%%%%%%%%%%%%%%%%%%%%%%%

Under the same assumptions,

\[
\boxed{
\lim_{x\to a}
[f(x)-g(x)]
=
L-M.
}
\]

%%%%%%%%%%%%%%%%%%%%%%%%%%%%%%%%%%%%%%%%%%%%%%%%%%%%%%%%%%%%
\subsection{Product Law}
\label{subsec:product-limit-law}
%%%%%%%%%%%%%%%%%%%%%%%%%%%%%%%%%%%%%%%%%%%%%%%%%%%%%%%%%%%%

If both limits exist,

\[
\boxed{
\lim_{x\to a}
f(x)g(x)
=
LM.
}
\]

%%%%%%%%%%%%%%%%%%%%%%%%%%%%%%%%%%%%%%%%%%%%%%%%%%%%%%%%%%%%
\subsection{Quotient Law}
\label{subsec:quotient-limit-law}
%%%%%%%%%%%%%%%%%%%%%%%%%%%%%%%%%%%%%%%%%%%%%%%%%%%%%%%%%%%%

If

\[
\lim_{x\to a}f(x)=L,
\]

\[
\lim_{x\to a}g(x)=M,
\]

and

\[
M\neq0,
\]

then

\[
\boxed{
\lim_{x\to a}
\frac{f(x)}{g(x)}
=
\frac{L}{M}.
}
\]

%%%%%%%%%%%%%%%%%%%%%%%%%%%%%%%%%%%%%%%%%%%%%%%%%%%%%%%%%%%%
\subsection{Power Law}
\label{subsec:power-limit-law}
%%%%%%%%%%%%%%%%%%%%%%%%%%%%%%%%%%%%%%%%%%%%%%%%%%%%%%%%%%%%

For positive integer \(n\),

\[
\boxed{
\lim_{x\to a}
[f(x)]^n
=
L^n
}
\]

whenever

\[
\lim_{x\to a}f(x)=L.
\]

%%%%%%%%%%%%%%%%%%%%%%%%%%%%%%%%%%%%%%%%%%%%%%%%%%%%%%%%%%%%
\subsection{Root Law}
\label{subsec:root-limit-law}
%%%%%%%%%%%%%%%%%%%%%%%%%%%%%%%%%%%%%%%%%%%%%%%%%%%%%%%%%%%%

For suitable domains,

\[
\boxed{
\lim_{x\to a}
\sqrt[n]{f(x)}
=
\sqrt[n]{L}.
}
\]

If \(n\) is even, the relevant quantities must remain in the real domain.

%%%%%%%%%%%%%%%%%%%%%%%%%%%%%%%%%%%%%%%%%%%%%%%%%%%%%%%%%%%%
\subsection{Absolute Value Law}
\label{subsec:absolute-value-limit-law}
%%%%%%%%%%%%%%%%%%%%%%%%%%%%%%%%%%%%%%%%%%%%%%%%%%%%%%%%%%%%

If

\[
\lim_{x\to a}f(x)=L,
\]

then

\[
\boxed{
\lim_{x\to a}|f(x)|
=
|L|.
}
\]

%%%%%%%%%%%%%%%%%%%%%%%%%%%%%%%%%%%%%%%%%%%%%%%%%%%%%%%%%%%%
\subsection{Polynomial Limits}
\label{subsec:polynomial-limit-reference}
%%%%%%%%%%%%%%%%%%%%%%%%%%%%%%%%%%%%%%%%%%%%%%%%%%%%%%%%%%%%

For a polynomial

\[
p(x),
\]

\[
\boxed{
\lim_{x\to a}p(x)
=
p(a).
}
\]

Thus polynomial limits can be evaluated by direct substitution.

%%%%%%%%%%%%%%%%%%%%%%%%%%%%%%%%%%%%%%%%%%%%%%%%%%%%%%%%%%%%
\subsection{Worked Example: Polynomial Limit}
\label{subsec:worked-polynomial-limit}
%%%%%%%%%%%%%%%%%%%%%%%%%%%%%%%%%%%%%%%%%%%%%%%%%%%%%%%%%%%%

\begin{workedexamplebox}{Direct Substitution}

Evaluate

\[
\lim_{x\to2}
(3x^2-5x+4).
\]

Since polynomials are continuous,

\[
3(2)^2-5(2)+4
=
12-10+4
=
6.
\]

\begin{answerbox}
\[
\boxed{
\lim_{x\to2}
(3x^2-5x+4)
=
6.
}
\]
\end{answerbox}

\end{workedexamplebox}

%%%%%%%%%%%%%%%%%%%%%%%%%%%%%%%%%%%%%%%%%%%%%%%%%%%%%%%%%%%%
\subsection{Rational Function Limits}
\label{subsec:rational-function-limit-reference}
%%%%%%%%%%%%%%%%%%%%%%%%%%%%%%%%%%%%%%%%%%%%%%%%%%%%%%%%%%%%

For rational function

\[
R(x)
=
\frac{p(x)}{q(x)},
\]

if

\[
q(a)\neq0,
\]

then

\[
\boxed{
\lim_{x\to a}
\frac{p(x)}{q(x)}
=
\frac{p(a)}{q(a)}.
}
\]

%%%%%%%%%%%%%%%%%%%%%%%%%%%%%%%%%%%%%%%%%%%%%%%%%%%%%%%%%%%%
\subsection{Direct Substitution}
\label{subsec:direct-substitution-limit}
%%%%%%%%%%%%%%%%%%%%%%%%%%%%%%%%%%%%%%%%%%%%%%%%%%%%%%%%%%%%

Direct substitution is appropriate when the function is continuous at the
point of interest.

For example,

\[
\boxed{
\lim_{x\to a}f(x)
=
f(a)
}
\]

whenever \(f\) is continuous at \(a\).

%%%%%%%%%%%%%%%%%%%%%%%%%%%%%%%%%%%%%%%%%%%%%%%%%%%%%%%%%%%%
\subsection{Indeterminate Forms}
\label{subsec:indeterminate-forms-reference}
%%%%%%%%%%%%%%%%%%%%%%%%%%%%%%%%%%%%%%%%%%%%%%%%%%%%%%%%%%%%

Expressions such as

\[
\boxed{
\frac00,
\qquad
\frac{\infty}{\infty},
\qquad
0\cdot\infty,
}
\]

\[
\boxed{
\infty-\infty,
\qquad
0^0,
\qquad
1^\infty,
\qquad
\infty^0
}
\]

are called indeterminate forms.

They are not answers.

They indicate that further analysis is required.

%%%%%%%%%%%%%%%%%%%%%%%%%%%%%%%%%%%%%%%%%%%%%%%%%%%%%%%%%%%%
\subsection{The Form \(0/0\)}
\label{subsec:zero-over-zero}
%%%%%%%%%%%%%%%%%%%%%%%%%%%%%%%%%%%%%%%%%%%%%%%%%%%%%%%%%%%%

If direct substitution gives

\[
\frac00,
\]

possible methods include:

\begin{itemize}
    \item factoring;
    \item rationalizing;
    \item simplifying algebraically;
    \item trigonometric identities;
    \item fundamental limits;
    \item L'Hospital's Rule when its hypotheses are satisfied;
    \item series expansions.
\end{itemize}

%%%%%%%%%%%%%%%%%%%%%%%%%%%%%%%%%%%%%%%%%%%%%%%%%%%%%%%%%%%%
\subsection{Factoring Method}
\label{subsec:factoring-limit-method}
%%%%%%%%%%%%%%%%%%%%%%%%%%%%%%%%%%%%%%%%%%%%%%%%%%%%%%%%%%%%

A removable factor can often be canceled before evaluating the limit.

For example,

\[
\frac{x^2-a^2}{x-a}
=
x+a
\]

for

\[
x\neq a.
\]

Thus

\[
\boxed{
\lim_{x\to a}
\frac{x^2-a^2}{x-a}
=
2a.
}
\]

%%%%%%%%%%%%%%%%%%%%%%%%%%%%%%%%%%%%%%%%%%%%%%%%%%%%%%%%%%%%
\subsection{Worked Example: Factoring}
\label{subsec:worked-factoring-limit}
%%%%%%%%%%%%%%%%%%%%%%%%%%%%%%%%%%%%%%%%%%%%%%%%%%%%%%%%%%%%

\begin{workedexamplebox}{Resolve \(0/0\) by Factoring}

Evaluate

\[
\lim_{x\to3}
\frac{x^2-9}{x-3}.
\]

Factor:

\[
x^2-9
=
(x-3)(x+3).
\]

For \(x\neq3\),

\[
\frac{x^2-9}{x-3}
=
x+3.
\]

Therefore,

\begin{answerbox}
\[
\boxed{
\lim_{x\to3}
\frac{x^2-9}{x-3}
=
6.
}
\]
\end{answerbox}

\end{workedexamplebox}

%%%%%%%%%%%%%%%%%%%%%%%%%%%%%%%%%%%%%%%%%%%%%%%%%%%%%%%%%%%%
\subsection{Rationalization Method}
\label{subsec:rationalization-limit-method}
%%%%%%%%%%%%%%%%%%%%%%%%%%%%%%%%%%%%%%%%%%%%%%%%%%%%%%%%%%%%

When radicals produce \(0/0\), multiply by a conjugate.

For example,

\[
\frac{
\sqrt{x}-\sqrt a
}{
x-a
}
\]

can be multiplied by

\[
\frac{
\sqrt{x}+\sqrt a
}{
\sqrt{x}+\sqrt a
}.
\]

Then

\[
\boxed{
\frac{
\sqrt{x}-\sqrt a
}{
x-a
}
=
\frac{
1
}{
\sqrt{x}+\sqrt a
}
}
\]

for \(x\neq a\).

%%%%%%%%%%%%%%%%%%%%%%%%%%%%%%%%%%%%%%%%%%%%%%%%%%%%%%%%%%%%
\subsection{Worked Example: Radical Limit}
\label{subsec:worked-radical-limit}
%%%%%%%%%%%%%%%%%%%%%%%%%%%%%%%%%%%%%%%%%%%%%%%%%%%%%%%%%%%%

\begin{workedexamplebox}{Use a Conjugate}

Evaluate

\[
\lim_{x\to4}
\frac{
\sqrt{x}-2
}{
x-4
}.
\]

Multiply by the conjugate:

\[
\frac{
\sqrt{x}-2
}{
x-4
}
\frac{
\sqrt{x}+2
}{
\sqrt{x}+2
}.
\]

Then

\[
\frac{
x-4
}{
(x-4)(\sqrt{x}+2)
}
=
\frac{
1
}{
\sqrt{x}+2
}.
\]

Therefore,

\begin{answerbox}
\[
\boxed{
\lim_{x\to4}
\frac{
\sqrt{x}-2
}{
x-4
}
=
\frac14.
}
\]
\end{answerbox}

\end{workedexamplebox}

%%%%%%%%%%%%%%%%%%%%%%%%%%%%%%%%%%%%%%%%%%%%%%%%%%%%%%%%%%%%
\subsection{Piecewise-Function Limits}
\label{subsec:piecewise-limit-reference}
%%%%%%%%%%%%%%%%%%%%%%%%%%%%%%%%%%%%%%%%%%%%%%%%%%%%%%%%%%%%

For a piecewise function, evaluate left- and right-hand limits separately.

For example,

\[
f(x)
=
\begin{cases}
x+1, & x<2,\\
x^2, & x\geq2.
\end{cases}
\]

Then

\[
\lim_{x\to2^-}f(x)=3,
\]

while

\[
\lim_{x\to2^+}f(x)=4.
\]

Therefore,

\[
\boxed{
\lim_{x\to2}f(x)
\text{ does not exist}.
}
\]

%%%%%%%%%%%%%%%%%%%%%%%%%%%%%%%%%%%%%%%%%%%%%%%%%%%%%%%%%%%%
\subsection{The Squeeze Theorem}
\label{subsec:squeeze-theorem-reference}
%%%%%%%%%%%%%%%%%%%%%%%%%%%%%%%%%%%%%%%%%%%%%%%%%%%%%%%%%%%%

Suppose

\[
g(x)
\leq
f(x)
\leq
h(x)
\]

near \(a\), and

\[
\lim_{x\to a}g(x)
=
\lim_{x\to a}h(x)
=
L.
\]

Then

\[
\boxed{
\lim_{x\to a}f(x)
=
L.
}
\]

%%%%%%%%%%%%%%%%%%%%%%%%%%%%%%%%%%%%%%%%%%%%%%%%%%%%%%%%%%%%
\subsection{Worked Example: Squeeze Theorem}
\label{subsec:worked-squeeze-theorem}
%%%%%%%%%%%%%%%%%%%%%%%%%%%%%%%%%%%%%%%%%%%%%%%%%%%%%%%%%%%%

\begin{workedexamplebox}{Oscillatory Function}

Evaluate

\[
\lim_{x\to0}
x^2\sin\left(\frac1x\right).
\]

Since

\[
-1
\leq
\sin\left(\frac1x\right)
\leq
1,
\]

multiplying by \(x^2\geq0\) gives

\[
-x^2
\leq
x^2\sin\left(\frac1x\right)
\leq
x^2.
\]

Since

\[
\lim_{x\to0}(-x^2)=0
\]

and

\[
\lim_{x\to0}x^2=0,
\]

the squeeze theorem gives

\begin{answerbox}
\[
\boxed{
\lim_{x\to0}
x^2\sin\left(\frac1x\right)
=
0.
}
\]
\end{answerbox}

\end{workedexamplebox}

%%%%%%%%%%%%%%%%%%%%%%%%%%%%%%%%%%%%%%%%%%%%%%%%%%%%%%%%%%%%
\subsection{Fundamental Sine Limit}
\label{subsec:fundamental-sine-limit}
%%%%%%%%%%%%%%%%%%%%%%%%%%%%%%%%%%%%%%%%%%%%%%%%%%%%%%%%%%%%

When angles are measured in radians,

\[
\boxed{
\lim_{x\to0}
\frac{\sin x}{x}
=
1.
}
\]

This is one of the fundamental limits of calculus.

%%%%%%%%%%%%%%%%%%%%%%%%%%%%%%%%%%%%%%%%%%%%%%%%%%%%%%%%%%%%
\subsection{Related Sine Limits}
\label{subsec:related-sine-limits}
%%%%%%%%%%%%%%%%%%%%%%%%%%%%%%%%%%%%%%%%%%%%%%%%%%%%%%%%%%%%

For constant \(a\),

\[
\boxed{
\lim_{x\to0}
\frac{\sin(ax)}{x}
=
a.
}
\]

Also,

\[
\boxed{
\lim_{x\to0}
\frac{\sin(ax)}{bx}
=
\frac{a}{b}
}
\]

when \(b\neq0\).

%%%%%%%%%%%%%%%%%%%%%%%%%%%%%%%%%%%%%%%%%%%%%%%%%%%%%%%%%%%%
\subsection{Worked Example: Trigonometric Limit}
\label{subsec:worked-trigonometric-limit}
%%%%%%%%%%%%%%%%%%%%%%%%%%%%%%%%%%%%%%%%%%%%%%%%%%%%%%%%%%%%

\begin{workedexamplebox}{Scaled Sine Limit}

Evaluate

\[
\lim_{x\to0}
\frac{\sin5x}{x}.
\]

Write

\[
\frac{\sin5x}{x}
=
5
\frac{\sin5x}{5x}.
\]

Then

\[
\lim_{x\to0}
\frac{\sin5x}{5x}
=
1.
\]

Therefore,

\begin{answerbox}
\[
\boxed{
\lim_{x\to0}
\frac{\sin5x}{x}
=
5.
}
\]
\end{answerbox}

\end{workedexamplebox}

%%%%%%%%%%%%%%%%%%%%%%%%%%%%%%%%%%%%%%%%%%%%%%%%%%%%%%%%%%%%
\subsection{Tangent Fundamental Limit}
\label{subsec:tangent-fundamental-limit}
%%%%%%%%%%%%%%%%%%%%%%%%%%%%%%%%%%%%%%%%%%%%%%%%%%%%%%%%%%%%

Since

\[
\tan x
=
\frac{\sin x}{\cos x},
\]

\[
\boxed{
\lim_{x\to0}
\frac{\tan x}{x}
=
1.
}
\]

%%%%%%%%%%%%%%%%%%%%%%%%%%%%%%%%%%%%%%%%%%%%%%%%%%%%%%%%%%%%
\subsection{Cosine Fundamental Limit}
\label{subsec:cosine-fundamental-limit}
%%%%%%%%%%%%%%%%%%%%%%%%%%%%%%%%%%%%%%%%%%%%%%%%%%%%%%%%%%%%

A standard result is

\[
\boxed{
\lim_{x\to0}
\frac{1-\cos x}{x^2}
=
\frac12.
}
\]

Equivalently,

\[
\boxed{
1-\cos x
\sim
\frac{x^2}{2}
}
\]

as \(x\to0\).

%%%%%%%%%%%%%%%%%%%%%%%%%%%%%%%%%%%%%%%%%%%%%%%%%%%%%%%%%%%%
\subsection{Worked Example: Cosine Limit}
\label{subsec:worked-cosine-limit}
%%%%%%%%%%%%%%%%%%%%%%%%%%%%%%%%%%%%%%%%%%%%%%%%%%%%%%%%%%%%

\begin{workedexamplebox}{Scaled Cosine Limit}

Evaluate

\[
\lim_{x\to0}
\frac{
1-\cos3x
}{
x^2
}.
\]

Write

\[
\frac{
1-\cos3x
}{
x^2
}
=
9
\frac{
1-\cos3x
}{
(3x)^2
}.
\]

Therefore,

\[
\begin{answerbox}
\[
\boxed{
\lim_{x\to0}
\frac{
1-\cos3x
}{
x^2
}
=
\frac92.
}
\]
\end{answerbox}

\end{workedexamplebox}

%%%%%%%%%%%%%%%%%%%%%%%%%%%%%%%%%%%%%%%%%%%%%%%%%%%%%%%%%%%%
\subsection{Exponential Fundamental Limit}
\label{subsec:exponential-fundamental-limit}
%%%%%%%%%%%%%%%%%%%%%%%%%%%%%%%%%%%%%%%%%%%%%%%%%%%%%%%%%%%%

A standard limit is

\[
\boxed{
\lim_{x\to0}
\frac{e^x-1}{x}
=
1.
}
\]

More generally,

\[
\boxed{
\lim_{x\to0}
\frac{e^{ax}-1}{x}
=
a.
}
\]

%%%%%%%%%%%%%%%%%%%%%%%%%%%%%%%%%%%%%%%%%%%%%%%%%%%%%%%%%%%%
\subsection{Logarithmic Fundamental Limit}
\label{subsec:logarithmic-fundamental-limit}
%%%%%%%%%%%%%%%%%%%%%%%%%%%%%%%%%%%%%%%%%%%%%%%%%%%%%%%%%%%%

A related logarithmic limit is

\[
\boxed{
\lim_{x\to0}
\frac{\ln(1+x)}{x}
=
1.
}
\]

More generally,

\[
\boxed{
\ln(1+x)
\sim
x
}
\]

as \(x\to0\).

%%%%%%%%%%%%%%%%%%%%%%%%%%%%%%%%%%%%%%%%%%%%%%%%%%%%%%%%%%%%
\subsection{Definition of \(e\) by a Limit}
\label{subsec:e-limit-reference}
%%%%%%%%%%%%%%%%%%%%%%%%%%%%%%%%%%%%%%%%%%%%%%%%%%%%%%%%%%%%

One defining limit is

\[
\boxed{
e
=
\lim_{n\to\infty}
\left(
1+\frac1n
\right)^n.
}
\]

A continuous-variable version is

\[
\boxed{
\lim_{x\to0}
(1+x)^{1/x}
=
e.
}
\]

%%%%%%%%%%%%%%%%%%%%%%%%%%%%%%%%%%%%%%%%%%%%%%%%%%%%%%%%%%%%
\subsection{General Exponential-Type Limit}
\label{subsec:general-exponential-type-limit}
%%%%%%%%%%%%%%%%%%%%%%%%%%%%%%%%%%%%%%%%%%%%%%%%%%%%%%%%%%%%

For suitable real \(a\),

\[
\boxed{
\lim_{x\to0}
(1+ax)^{1/x}
=
e^a.
}
\]

Similarly,

\[
\boxed{
\lim_{n\to\infty}
\left(
1+\frac{a}{n}
\right)^n
=
e^a.
}
\]

%%%%%%%%%%%%%%%%%%%%%%%%%%%%%%%%%%%%%%%%%%%%%%%%%%%%%%%%%%%%
\subsection{Composition Law for Limits}
\label{subsec:composition-limit-law}
%%%%%%%%%%%%%%%%%%%%%%%%%%%%%%%%%%%%%%%%%%%%%%%%%%%%%%%%%%%%

Suppose

\[
\lim_{x\to a}g(x)=L
\]

and \(f\) is continuous at \(L\).

Then

\[
\boxed{
\lim_{x\to a}
f(g(x))
=
f(L).
}
\]

This justifies direct substitution into continuous outer functions.

%%%%%%%%%%%%%%%%%%%%%%%%%%%%%%%%%%%%%%%%%%%%%%%%%%%%%%%%%%%%
\subsection{Continuity and Limits}
\label{subsec:continuity-limit-reference}
%%%%%%%%%%%%%%%%%%%%%%%%%%%%%%%%%%%%%%%%%%%%%%%%%%%%%%%%%%%%

A function \(f\) is continuous at \(a\) when

\[
\boxed{
\lim_{x\to a}f(x)
=
f(a).
}
\]

This requires:

\begin{enumerate}
    \item \(f(a)\) is defined;
    \item \(\lim_{x\to a}f(x)\) exists;
    \item the limit equals the function value.
\end{enumerate}

%%%%%%%%%%%%%%%%%%%%%%%%%%%%%%%%%%%%%%%%%%%%%%%%%%%%%%%%%%%%
\subsection{Continuous Functions Commonly Used}
\label{subsec:common-continuous-functions-reference}
%%%%%%%%%%%%%%%%%%%%%%%%%%%%%%%%%%%%%%%%%%%%%%%%%%%%%%%%%%%%

The following are continuous on their natural domains:

\begin{itemize}
    \item polynomials;
    \item rational functions where the denominator is nonzero;
    \item exponential functions;
    \item logarithmic functions on positive arguments;
    \item sine and cosine;
    \item tangent where cosine is nonzero;
    \item root functions on their real domains;
    \item inverse trigonometric functions on their domains.
\end{itemize}

%%%%%%%%%%%%%%%%%%%%%%%%%%%%%%%%%%%%%%%%%%%%%%%%%%%%%%%%%%%%
\subsection{Limits at Infinity}
\label{subsec:limits-at-infinity-reference}
%%%%%%%%%%%%%%%%%%%%%%%%%%%%%%%%%%%%%%%%%%%%%%%%%%%%%%%%%%%%

The notation

\[
\boxed{
\lim_{x\to\infty}f(x)=L
}
\]

describes the behavior of \(f(x)\) as \(x\) grows without bound.

Similarly,

\[
\boxed{
\lim_{x\to-\infty}f(x)=L
}
\]

describes behavior as \(x\) becomes arbitrarily negative.

%%%%%%%%%%%%%%%%%%%%%%%%%%%%%%%%%%%%%%%%%%%%%%%%%%%%%%%%%%%%
\subsection{Horizontal Asymptotes}
\label{subsec:horizontal-asymptote-limit}
%%%%%%%%%%%%%%%%%%%%%%%%%%%%%%%%%%%%%%%%%%%%%%%%%%%%%%%%%%%%

If

\[
\boxed{
\lim_{x\to\infty}f(x)=L
}
\]

or

\[
\boxed{
\lim_{x\to-\infty}f(x)=L,
}
\]

then

\[
\boxed{
y=L
}
\]

is a horizontal asymptote in the corresponding direction.

%%%%%%%%%%%%%%%%%%%%%%%%%%%%%%%%%%%%%%%%%%%%%%%%%%%%%%%%%%%%
\subsection{Basic Reciprocal Limits}
\label{subsec:basic-reciprocal-infinity-limits}
%%%%%%%%%%%%%%%%%%%%%%%%%%%%%%%%%%%%%%%%%%%%%%%%%%%%%%%%%%%%

For positive integer \(n\),

\[
\boxed{
\lim_{x\to\infty}
\frac{1}{x^n}
=
0.
}
\]

Similarly,

\[
\boxed{
\lim_{x\to-\infty}
\frac{1}{x^n}
=
0.
}
\]

%%%%%%%%%%%%%%%%%%%%%%%%%%%%%%%%%%%%%%%%%%%%%%%%%%%%%%%%%%%%
\subsection{Polynomial Dominant Terms}
\label{subsec:polynomial-dominant-term}
%%%%%%%%%%%%%%%%%%%%%%%%%%%%%%%%%%%%%%%%%%%%%%%%%%%%%%%%%%%%

For polynomial

\[
p(x)
=
a_nx^n+\cdots+a_0,
\]

the highest-degree term determines the dominant growth as

\[
|x|\to\infty.
\]

Thus

\[
\boxed{
p(x)
\sim
a_nx^n.
}
\]

%%%%%%%%%%%%%%%%%%%%%%%%%%%%%%%%%%%%%%%%%%%%%%%%%%%%%%%%%%%%
\subsection{Worked Example: Polynomial at Infinity}
\label{subsec:worked-polynomial-infinity}
%%%%%%%%%%%%%%%%%%%%%%%%%%%%%%%%%%%%%%%%%%%%%%%%%%%%%%%%%%%%

\begin{workedexamplebox}{Dominant Term}

Evaluate

\[
\lim_{x\to\infty}
(3x^4-2x+7).
\]

The dominant term is

\[
3x^4.
\]

Since

\[
3x^4\to\infty,
\]

we have

\begin{answerbox}
\[
\boxed{
\lim_{x\to\infty}
(3x^4-2x+7)
=
\infty.
}
\]
\end{answerbox}

\end{workedexamplebox}

%%%%%%%%%%%%%%%%%%%%%%%%%%%%%%%%%%%%%%%%%%%%%%%%%%%%%%%%%%%%
\subsection{Rational Functions at Infinity}
\label{subsec:rational-functions-infinity}
%%%%%%%%%%%%%%%%%%%%%%%%%%%%%%%%%%%%%%%%%%%%%%%%%%%%%%%%%%%%

Consider

\[
\frac{
a_nx^n+\cdots
}{
b_mx^m+\cdots
}.
\]

The behavior depends on the degrees \(n\) and \(m\).

%%%%%%%%%%%%%%%%%%%%%%%%%%%%%%%%%%%%%%%%%%%%%%%%%%%%%%%%%%%%
\subsection{Numerator Degree Less Than Denominator}
\label{subsec:rational-degree-less}
%%%%%%%%%%%%%%%%%%%%%%%%%%%%%%%%%%%%%%%%%%%%%%%%%%%%%%%%%%%%

If

\[
n<m,
\]

then

\[
\boxed{
\lim_{x\to\pm\infty}
\frac{p(x)}{q(x)}
=
0.
}
\]

%%%%%%%%%%%%%%%%%%%%%%%%%%%%%%%%%%%%%%%%%%%%%%%%%%%%%%%%%%%%
\subsection{Equal Degrees}
\label{subsec:rational-equal-degree}
%%%%%%%%%%%%%%%%%%%%%%%%%%%%%%%%%%%%%%%%%%%%%%%%%%%%%%%%%%%%

If

\[
n=m,
\]

then

\[
\boxed{
\lim_{x\to\pm\infty}
\frac{p(x)}{q(x)}
=
\frac{a_n}{b_n}.
}
\]

Thus the limit is the ratio of leading coefficients.

%%%%%%%%%%%%%%%%%%%%%%%%%%%%%%%%%%%%%%%%%%%%%%%%%%%%%%%%%%%%
\subsection{Numerator Degree Greater Than Denominator}
\label{subsec:rational-degree-greater}
%%%%%%%%%%%%%%%%%%%%%%%%%%%%%%%%%%%%%%%%%%%%%%%%%%%%%%%%%%%%

If

\[
n>m,
\]

then the rational function generally grows without bound in magnitude or
behaves like a polynomial of degree \(n-m\).

Polynomial long division can reveal the asymptotic behavior.

%%%%%%%%%%%%%%%%%%%%%%%%%%%%%%%%%%%%%%%%%%%%%%%%%%%%%%%%%%%%
\subsection{Worked Example: Rational Limit at Infinity}
\label{subsec:worked-rational-infinity}
%%%%%%%%%%%%%%%%%%%%%%%%%%%%%%%%%%%%%%%%%%%%%%%%%%%%%%%%%%%%

\begin{workedexamplebox}{Equal Degrees}

Evaluate

\[
\lim_{x\to\infty}
\frac{
5x^3-2x
}{
2x^3+7
}.
\]

Divide numerator and denominator by \(x^3\):

\[
\frac{
5-\frac{2}{x^2}
}{
2+\frac{7}{x^3}
}.
\]

As \(x\to\infty\),

\[
\frac{2}{x^2}\to0,
\qquad
\frac{7}{x^3}\to0.
\]

Therefore,

\begin{answerbox}
\[
\boxed{
\lim_{x\to\infty}
\frac{
5x^3-2x
}{
2x^3+7
}
=
\frac52.
}
\]
\end{answerbox}

\end{workedexamplebox}

%%%%%%%%%%%%%%%%%%%%%%%%%%%%%%%%%%%%%%%%%%%%%%%%%%%%%%%%%%%%
\subsection{Radicals at Infinity}
\label{subsec:radicals-at-infinity}
%%%%%%%%%%%%%%%%%%%%%%%%%%%%%%%%%%%%%%%%%%%%%%%%%%%%%%%%%%%%

When factoring powers from radicals, remember:

\[
\boxed{
\sqrt{x^2}=|x|.
}
\]

For example,

\[
\sqrt{x^2+1}
=
|x|
\sqrt{
1+\frac1{x^2}
}.
\]

As \(x\to+\infty\),

\[
|x|=x.
\]

As \(x\to-\infty\),

\[
|x|=-x.
\]

%%%%%%%%%%%%%%%%%%%%%%%%%%%%%%%%%%%%%%%%%%%%%%%%%%%%%%%%%%%%
\subsection{Worked Example: Radical at Infinity}
\label{subsec:worked-radical-infinity}
%%%%%%%%%%%%%%%%%%%%%%%%%%%%%%%%%%%%%%%%%%%%%%%%%%%%%%%%%%%%

\begin{workedexamplebox}{Factor Correctly From a Square Root}

Evaluate

\[
\lim_{x\to-\infty}
\frac{
\sqrt{x^2+4}
}{
x
}.
\]

Factor \(x^2\) from the radical:

\[
\sqrt{x^2+4}
=
|x|
\sqrt{
1+\frac4{x^2}
}.
\]

For \(x\to-\infty\),

\[
|x|=-x.
\]

Therefore,

\[
\frac{
\sqrt{x^2+4}
}{
x
}
=
-
\sqrt{
1+\frac4{x^2}
}.
\]

Hence

\begin{answerbox}
\[
\boxed{
\lim_{x\to-\infty}
\frac{
\sqrt{x^2+4}
}{
x
}
=
-1.
}
\]
\end{answerbox}

\end{workedexamplebox}

%%%%%%%%%%%%%%%%%%%%%%%%%%%%%%%%%%%%%%%%%%%%%%%%%%%%%%%%%%%%
\subsection{Infinite Limits}
\label{subsec:infinite-limits-reference}
%%%%%%%%%%%%%%%%%%%%%%%%%%%%%%%%%%%%%%%%%%%%%%%%%%%%%%%%%%%%

The notation

\[
\boxed{
\lim_{x\to a}f(x)
=
\infty
}
\]

means \(f(x)\) grows arbitrarily large as \(x\to a\).

This does not mean that infinity is an ordinary real-number output.

%%%%%%%%%%%%%%%%%%%%%%%%%%%%%%%%%%%%%%%%%%%%%%%%%%%%%%%%%%%%
\subsection{Vertical Asymptotes}
\label{subsec:vertical-asymptote-limit}
%%%%%%%%%%%%%%%%%%%%%%%%%%%%%%%%%%%%%%%%%%%%%%%%%%%%%%%%%%%%

If

\[
\lim_{x\to a^-}f(x)
=
\pm\infty
\]

or

\[
\lim_{x\to a^+}f(x)
=
\pm\infty,
\]

then

\[
\boxed{
x=a
}
\]

is typically a vertical asymptote.

%%%%%%%%%%%%%%%%%%%%%%%%%%%%%%%%%%%%%%%%%%%%%%%%%%%%%%%%%%%%
\subsection{Basic Reciprocal Infinite Limit}
\label{subsec:reciprocal-infinite-limit}
%%%%%%%%%%%%%%%%%%%%%%%%%%%%%%%%%%%%%%%%%%%%%%%%%%%%%%%%%%%%

For example,

\[
\boxed{
\lim_{x\to0^+}
\frac1x
=
+\infty,
}
\]

while

\[
\boxed{
\lim_{x\to0^-}
\frac1x
=
-\infty.
}
\]

Thus the two-sided limit

\[
\lim_{x\to0}\frac1x
\]

does not exist as a single infinite limit.

%%%%%%%%%%%%%%%%%%%%%%%%%%%%%%%%%%%%%%%%%%%%%%%%%%%%%%%%%%%%
\subsection{Even Reciprocal Powers}
\label{subsec:even-reciprocal-infinite-limit}
%%%%%%%%%%%%%%%%%%%%%%%%%%%%%%%%%%%%%%%%%%%%%%%%%%%%%%%%%%%%

For even positive integer \(n\),

\[
\boxed{
\lim_{x\to0}
\frac1{x^n}
=
+\infty.
}
\]

Both sides approach positive infinity.

%%%%%%%%%%%%%%%%%%%%%%%%%%%%%%%%%%%%%%%%%%%%%%%%%%%%%%%%%%%%
\subsection{Odd Reciprocal Powers}
\label{subsec:odd-reciprocal-infinite-limit}
%%%%%%%%%%%%%%%%%%%%%%%%%%%%%%%%%%%%%%%%%%%%%%%%%%%%%%%%%%%%

For odd positive integer \(n\),

\[
\boxed{
\lim_{x\to0^+}
\frac1{x^n}
=
+\infty,
}
\]

while

\[
\boxed{
\lim_{x\to0^-}
\frac1{x^n}
=
-\infty.
}
\]

%%%%%%%%%%%%%%%%%%%%%%%%%%%%%%%%%%%%%%%%%%%%%%%%%%%%%%%%%%%%
\subsection{Infinity Is Not a Number in Ordinary Limit Algebra}
\label{subsec:infinity-not-number}
%%%%%%%%%%%%%%%%%%%%%%%%%%%%%%%%%%%%%%%%%%%%%%%%%%%%%%%%%%%%

One should not manipulate infinity as though it were an ordinary real number.

Expressions such as

\[
\infty-\infty
\]

or

\[
\frac{\infty}{\infty}
\]

are not defined by ordinary arithmetic.

They are indeterminate forms describing competing growth.

%%%%%%%%%%%%%%%%%%%%%%%%%%%%%%%%%%%%%%%%%%%%%%%%%%%%%%%%%%%%
\subsection{Growth Hierarchy}
\label{subsec:growth-hierarchy-reference}
%%%%%%%%%%%%%%%%%%%%%%%%%%%%%%%%%%%%%%%%%%%%%%%%%%%%%%%%%%%%

As \(x\to\infty\), a standard growth hierarchy is

\[
\boxed{
\log x
\ll
x^\alpha
\ll
a^x
}
\]

for

\[
\alpha>0,
\qquad
a>1.
\]

For sequences, one often extends this to

\[
\boxed{
\log n
\ll
n^\alpha
\ll
a^n
\ll
n!
\ll
n^n.
}
\]

%%%%%%%%%%%%%%%%%%%%%%%%%%%%%%%%%%%%%%%%%%%%%%%%%%%%%%%%%%%%
\subsection{Logarithms Versus Powers}
\label{subsec:logarithm-versus-power}
%%%%%%%%%%%%%%%%%%%%%%%%%%%%%%%%%%%%%%%%%%%%%%%%%%%%%%%%%%%%

For every \(\alpha>0\),

\[
\boxed{
\lim_{x\to\infty}
\frac{\ln x}{x^\alpha}
=
0.
}
\]

Thus every positive power eventually grows faster than the logarithm.

%%%%%%%%%%%%%%%%%%%%%%%%%%%%%%%%%%%%%%%%%%%%%%%%%%%%%%%%%%%%
\subsection{Powers Versus Exponentials}
\label{subsec:power-versus-exponential}
%%%%%%%%%%%%%%%%%%%%%%%%%%%%%%%%%%%%%%%%%%%%%%%%%%%%%%%%%%%%

For \(a>1\) and every positive integer \(n\),

\[
\boxed{
\lim_{x\to\infty}
\frac{x^n}{a^x}
=
0.
}
\]

Thus exponential growth eventually dominates polynomial growth.

%%%%%%%%%%%%%%%%%%%%%%%%%%%%%%%%%%%%%%%%%%%%%%%%%%%%%%%%%%%%
\subsection{Exponentials Versus Factorials}
\label{subsec:exponential-versus-factorial}
%%%%%%%%%%%%%%%%%%%%%%%%%%%%%%%%%%%%%%%%%%%%%%%%%%%%%%%%%%%%

For every fixed \(a>0\),

\[
\boxed{
\lim_{n\to\infty}
\frac{a^n}{n!}
=
0.
}
\]

Thus factorial growth dominates any fixed-base exponential.

%%%%%%%%%%%%%%%%%%%%%%%%%%%%%%%%%%%%%%%%%%%%%%%%%%%%%%%%%%%%
\subsection{Sequence Limits}
\label{subsec:sequence-limits-reference}
%%%%%%%%%%%%%%%%%%%%%%%%%%%%%%%%%%%%%%%%%%%%%%%%%%%%%%%%%%%%

A sequence

\[
(a_n)
\]

converges to \(L\) if

\[
\boxed{
\lim_{n\to\infty}a_n=L.
}
\]

Formally, for every

\[
\varepsilon>0,
\]

there exists \(N\) such that

\[
n\geq N
\]

implies

\[
\boxed{
|a_n-L|<\varepsilon.
}
\]

%%%%%%%%%%%%%%%%%%%%%%%%%%%%%%%%%%%%%%%%%%%%%%%%%%%%%%%%%%%%
\subsection{Basic Sequence Examples}
\label{subsec:basic-sequence-limit-examples}
%%%%%%%%%%%%%%%%%%%%%%%%%%%%%%%%%%%%%%%%%%%%%%%%%%%%%%%%%%%%

\[
\boxed{
\lim_{n\to\infty}
\frac1n
=
0,
}
\]

\[
\boxed{
\lim_{n\to\infty}
\frac1{n^p}
=
0
\qquad
(p>0),
}
\]

and if

\[
|r|<1,
\]

then

\[
\boxed{
\lim_{n\to\infty}
r^n
=
0.
}
\]

%%%%%%%%%%%%%%%%%%%%%%%%%%%%%%%%%%%%%%%%%%%%%%%%%%%%%%%%%%%%
\subsection{Geometric Sequence Behavior}
\label{subsec:geometric-sequence-behavior}
%%%%%%%%%%%%%%%%%%%%%%%%%%%%%%%%%%%%%%%%%%%%%%%%%%%%%%%%%%%%

For

\[
a_n=r^n,
\]

the behavior depends on \(r\).

If

\[
|r|<1,
\]

then

\[
a_n\to0.
\]

If

\[
r=1,
\]

then

\[
a_n=1.
\]

If

\[
r>1,
\]

then

\[
a_n\to\infty.
\]

If

\[
r=-1,
\]

the sequence oscillates and does not converge.

%%%%%%%%%%%%%%%%%%%%%%%%%%%%%%%%%%%%%%%%%%%%%%%%%%%%%%%%%%%%
\subsection{Monotone Convergence Principle}
\label{subsec:monotone-convergence-sequence}
%%%%%%%%%%%%%%%%%%%%%%%%%%%%%%%%%%%%%%%%%%%%%%%%%%%%%%%%%%%%

A fundamental result for real sequences is:

\begin{importantbox}
Every monotone increasing sequence that is bounded above converges, and every
monotone decreasing sequence that is bounded below converges.
\end{importantbox}

This is a consequence of completeness of \(\mathbb{R}\).

%%%%%%%%%%%%%%%%%%%%%%%%%%%%%%%%%%%%%%%%%%%%%%%%%%%%%%%%%%%%
\subsection{Boundedness Alone Does Not Guarantee Convergence}
\label{subsec:bounded-not-convergent}
%%%%%%%%%%%%%%%%%%%%%%%%%%%%%%%%%%%%%%%%%%%%%%%%%%%%%%%%%%%%

The sequence

\[
\boxed{
a_n=(-1)^n
}
\]

is bounded because

\[
-1\leq a_n\leq1,
\]

but it does not converge.

Thus boundedness alone is insufficient.

%%%%%%%%%%%%%%%%%%%%%%%%%%%%%%%%%%%%%%%%%%%%%%%%%%%%%%%%%%%%
\subsection{Subsequence Principle}
\label{subsec:subsequence-limit-principle}
%%%%%%%%%%%%%%%%%%%%%%%%%%%%%%%%%%%%%%%%%%%%%%%%%%%%%%%%%%%%

If

\[
a_n\to L,
\]

then every subsequence also converges to \(L\).

Therefore, if two subsequences converge to different limits, the original
sequence cannot converge.

%%%%%%%%%%%%%%%%%%%%%%%%%%%%%%%%%%%%%%%%%%%%%%%%%%%%%%%%%%%%
\subsection{Worked Example: Oscillatory Sequence}
\label{subsec:worked-oscillatory-sequence}
%%%%%%%%%%%%%%%%%%%%%%%%%%%%%%%%%%%%%%%%%%%%%%%%%%%%%%%%%%%%

\begin{workedexamplebox}{Show a Sequence Diverges}

Consider

\[
a_n=(-1)^n.
\]

The even subsequence satisfies

\[
a_{2n}=1.
\]

The odd subsequence satisfies

\[
a_{2n+1}=-1.
\]

Since the two subsequences have different limits,

\begin{answerbox}
\[
\boxed{
(-1)^n
\text{ does not converge}.
}
\]
\end{answerbox}

\end{workedexamplebox}

%%%%%%%%%%%%%%%%%%%%%%%%%%%%%%%%%%%%%%%%%%%%%%%%%%%%%%%%%%%%
\subsection{Recursive Sequence Limits}
\label{subsec:recursive-sequence-limits}
%%%%%%%%%%%%%%%%%%%%%%%%%%%%%%%%%%%%%%%%%%%%%%%%%%%%%%%%%%%%

If a recursive sequence

\[
a_{n+1}=F(a_n)
\]

converges to \(L\), and \(F\) is continuous at \(L\), then

\[
\boxed{
L=F(L).
}
\]

This fixed-point equation gives candidate limits.

It does not by itself prove convergence.

%%%%%%%%%%%%%%%%%%%%%%%%%%%%%%%%%%%%%%%%%%%%%%%%%%%%%%%%%%%%
\subsection{Worked Example: Recursive Limit Candidate}
\label{subsec:worked-recursive-limit}
%%%%%%%%%%%%%%%%%%%%%%%%%%%%%%%%%%%%%%%%%%%%%%%%%%%%%%%%%%%%

\begin{workedexamplebox}{Find a Possible Limit}

Suppose

\[
a_{n+1}
=
\sqrt{2+a_n}
\]

and assume the sequence converges to \(L\geq0\).

Then

\[
L
=
\sqrt{2+L}.
\]

Squaring gives

\[
L^2-L-2=0.
\]

Thus

\[
(L-2)(L+1)=0.
\]

Since \(L\geq0\),

\begin{answerbox}
\[
\boxed{
L=2.
}
\]
\end{answerbox}

This calculation identifies the possible limit; a complete proof must also
establish convergence.

\end{workedexamplebox}

%%%%%%%%%%%%%%%%%%%%%%%%%%%%%%%%%%%%%%%%%%%%%%%%%%%%%%%%%%%%
\subsection{Asymptotic Equivalence}
\label{subsec:asymptotic-equivalence-limit}
%%%%%%%%%%%%%%%%%%%%%%%%%%%%%%%%%%%%%%%%%%%%%%%%%%%%%%%%%%%%

The notation

\[
\boxed{
f(x)\sim g(x)
}
\]

as \(x\to a\) means

\[
\boxed{
\frac{f(x)}{g(x)}
\to1.
}
\]

This says that \(f\) and \(g\) have the same leading-order behavior.

%%%%%%%%%%%%%%%%%%%%%%%%%%%%%%%%%%%%%%%%%%%%%%%%%%%%%%%%%%%%
\subsection{Common Local Equivalences}
\label{subsec:common-local-equivalences}
%%%%%%%%%%%%%%%%%%%%%%%%%%%%%%%%%%%%%%%%%%%%%%%%%%%%%%%%%%%%

As \(x\to0\),

\[
\boxed{
\sin x
\sim
x,
}
\]

\[
\boxed{
\tan x
\sim
x,
}
\]

\[
\boxed{
e^x-1
\sim
x,
}
\]

\[
\boxed{
\ln(1+x)
\sim
x,
}
\]

and

\[
\boxed{
1-\cos x
\sim
\frac{x^2}{2}.
}
\]

%%%%%%%%%%%%%%%%%%%%%%%%%%%%%%%%%%%%%%%%%%%%%%%%%%%%%%%%%%%%
\subsection{Using Asymptotic Equivalence}
\label{subsec:using-asymptotic-equivalence}
%%%%%%%%%%%%%%%%%%%%%%%%%%%%%%%%%%%%%%%%%%%%%%%%%%%%%%%%%%%%

For example,

\[
\lim_{x\to0}
\frac{
\sin3x
}{
\ln(1+5x)
}
\]

can be analyzed using

\[
\sin3x\sim3x
\]

and

\[
\ln(1+5x)\sim5x.
\]

Therefore,

\[
\boxed{
\lim_{x\to0}
\frac{
\sin3x
}{
\ln(1+5x)
}
=
\frac35.
}
\]

%%%%%%%%%%%%%%%%%%%%%%%%%%%%%%%%%%%%%%%%%%%%%%%%%%%%%%%%%%%%
\subsection{Big-O Notation}
\label{subsec:big-o-limit-reference}
%%%%%%%%%%%%%%%%%%%%%%%%%%%%%%%%%%%%%%%%%%%%%%%%%%%%%%%%%%%%

The notation

\[
\boxed{
f(x)=O(g(x))
}
\]

as \(x\to a\) means, roughly, that \(|f(x)|\) is bounded by a constant
multiple of \(|g(x)|\) near the limiting regime.

More precisely, there exists \(C>0\) such that

\[
\boxed{
|f(x)|
\leq
C|g(x)|
}
\]

sufficiently near the relevant limit.

%%%%%%%%%%%%%%%%%%%%%%%%%%%%%%%%%%%%%%%%%%%%%%%%%%%%%%%%%%%%
\subsection{Little-o Notation}
\label{subsec:little-o-limit-reference}
%%%%%%%%%%%%%%%%%%%%%%%%%%%%%%%%%%%%%%%%%%%%%%%%%%%%%%%%%%%%

The notation

\[
\boxed{
f(x)=o(g(x))
}
\]

means

\[
\boxed{
\frac{f(x)}{g(x)}
\to0.
}
\]

Thus \(f\) is asymptotically negligible relative to \(g\).

%%%%%%%%%%%%%%%%%%%%%%%%%%%%%%%%%%%%%%%%%%%%%%%%%%%%%%%%%%%%
\subsection{Taylor-Type Local Approximations}
\label{subsec:taylor-local-limit-reference}
%%%%%%%%%%%%%%%%%%%%%%%%%%%%%%%%%%%%%%%%%%%%%%%%%%%%%%%%%%%%

Near \(x=0\),

\[
\boxed{
e^x
=
1+x+\frac{x^2}{2}
+
O(x^3),
}
\]

\[
\boxed{
\sin x
=
x-\frac{x^3}{6}
+
O(x^5),
}
\]

\[
\boxed{
\cos x
=
1-\frac{x^2}{2}
+
O(x^4),
}
\]

and

\[
\boxed{
\ln(1+x)
=
x-\frac{x^2}{2}
+
O(x^3).
}
\]

These expansions can evaluate complicated limits efficiently.

%%%%%%%%%%%%%%%%%%%%%%%%%%%%%%%%%%%%%%%%%%%%%%%%%%%%%%%%%%%%
\subsection{Worked Example: Series-Based Limit}
\label{subsec:worked-series-limit}
%%%%%%%%%%%%%%%%%%%%%%%%%%%%%%%%%%%%%%%%%%%%%%%%%%%%%%%%%%%%

\begin{workedexamplebox}{Use Leading Terms}

Evaluate

\[
\lim_{x\to0}
\frac{
e^x-1-x
}{
x^2
}.
\]

Using

\[
e^x
=
1+x+\frac{x^2}{2}
+
O(x^3),
\]

we obtain

\[
e^x-1-x
=
\frac{x^2}{2}
+
O(x^3).
\]

Divide by \(x^2\):

\[
\frac{
e^x-1-x
}{
x^2
}
=
\frac12+O(x).
\]

Therefore,

\begin{answerbox}
\[
\boxed{
\lim_{x\to0}
\frac{
e^x-1-x
}{
x^2
}
=
\frac12.
}
\]
\end{answerbox}

\end{workedexamplebox}

%%%%%%%%%%%%%%%%%%%%%%%%%%%%%%%%%%%%%%%%%%%%%%%%%%%%%%%%%%%%
\subsection{L'Hospital's Rule}
\label{subsec:lhospital-reference}
%%%%%%%%%%%%%%%%%%%%%%%%%%%%%%%%%%%%%%%%%%%%%%%%%%%%%%%%%%%%

Suppose a quotient produces an indeterminate form

\[
\frac00
\]

or

\[
\frac{\pm\infty}{\pm\infty},
\]

and the hypotheses of L'Hospital's Rule are satisfied.

Then under appropriate conditions,

\[
\boxed{
\lim_{x\to a}
\frac{f(x)}{g(x)}
=
\lim_{x\to a}
\frac{f'(x)}{g'(x)}
}
\]

provided the derivative quotient limit exists or is infinite in the required
sense.

This is not a general rule for arbitrary quotients.

%%%%%%%%%%%%%%%%%%%%%%%%%%%%%%%%%%%%%%%%%%%%%%%%%%%%%%%%%%%%
\subsection{Worked Example: L'Hospital's Rule}
\label{subsec:worked-lhospital}
%%%%%%%%%%%%%%%%%%%%%%%%%%%%%%%%%%%%%%%%%%%%%%%%%%%%%%%%%%%%

\begin{workedexamplebox}{A \(0/0\) Limit}

Evaluate

\[
\lim_{x\to0}
\frac{
e^x-1
}{
x
}.
\]

Direct substitution gives

\[
\frac00.
\]

Differentiate numerator and denominator:

\[
\frac{d}{dx}(e^x-1)=e^x,
\]

\[
\frac{d}{dx}x=1.
\]

Thus

\[
\lim_{x\to0}
\frac{
e^x-1
}{
x
}
=
\lim_{x\to0}
e^x.
\]

Therefore,

\begin{answerbox}
\[
\boxed{
1.
}
\]
\end{answerbox}

\end{workedexamplebox}

%%%%%%%%%%%%%%%%%%%%%%%%%%%%%%%%%%%%%%%%%%%%%%%%%%%%%%%%%%%%
\subsection{Repeated L'Hospital Applications}
\label{subsec:repeated-lhospital}
%%%%%%%%%%%%%%%%%%%%%%%%%%%%%%%%%%%%%%%%%%%%%%%%%%%%%%%%%%%%

If the derivative quotient remains indeterminate and the hypotheses remain
valid, L'Hospital's Rule may sometimes be applied again.

For example,

\[
\frac{
1-\cos x
}{
x^2
}
\]

gives \(0/0\).

One application gives

\[
\frac{
\sin x
}{
2x
},
\]

which is still \(0/0\).

A second application gives

\[
\frac{
\cos x
}{
2
}.
\]

Therefore,

\[
\boxed{
\lim_{x\to0}
\frac{
1-\cos x
}{
x^2
}
=
\frac12.
}
\]

%%%%%%%%%%%%%%%%%%%%%%%%%%%%%%%%%%%%%%%%%%%%%%%%%%%%%%%%%%%%
\subsection{Transforming \(0\cdot\infty\)}
\label{subsec:zero-times-infinity}
%%%%%%%%%%%%%%%%%%%%%%%%%%%%%%%%%%%%%%%%%%%%%%%%%%%%%%%%%%%%

A product such as

\[
f(x)g(x)
\]

with

\[
f(x)\to0
\]

and

\[
g(x)\to\infty
\]

has indeterminate form

\[
0\cdot\infty.
\]

Rewrite as a quotient:

\[
\boxed{
f(x)g(x)
=
\frac{f(x)}{1/g(x)}
}
\]

or

\[
\boxed{
f(x)g(x)
=
\frac{g(x)}{1/f(x)}.
}
\]

%%%%%%%%%%%%%%%%%%%%%%%%%%%%%%%%%%%%%%%%%%%%%%%%%%%%%%%%%%%%
\subsection{Transforming \(\infty-\infty\)}
\label{subsec:infinity-minus-infinity}
%%%%%%%%%%%%%%%%%%%%%%%%%%%%%%%%%%%%%%%%%%%%%%%%%%%%%%%%%%%%

Expressions with form

\[
\infty-\infty
\]

often require:

\begin{itemize}
    \item a common denominator;
    \item factoring dominant terms;
    \item rationalization;
    \item logarithmic combination;
    \item algebraic simplification.
\end{itemize}

%%%%%%%%%%%%%%%%%%%%%%%%%%%%%%%%%%%%%%%%%%%%%%%%%%%%%%%%%%%%
\subsection{Worked Example: Infinity Minus Infinity}
\label{subsec:worked-infinity-minus-infinity}
%%%%%%%%%%%%%%%%%%%%%%%%%%%%%%%%%%%%%%%%%%%%%%%%%%%%%%%%%%%%

\begin{workedexamplebox}{Rationalize at Infinity}

Evaluate

\[
\lim_{x\to\infty}
\left(
\sqrt{x^2+x}-x
\right).
\]

Multiply by the conjugate:

\[
\sqrt{x^2+x}-x
=
\frac{
x
}{
\sqrt{x^2+x}+x
}.
\]

For \(x>0\),

\[
\sqrt{x^2+x}
=
x
\sqrt{
1+\frac1x
}.
\]

Therefore,

\[
\frac{x}{
x\sqrt{1+1/x}+x
}
=
\frac{
1
}{
\sqrt{1+1/x}+1
}.
\]

Hence

\begin{answerbox}
\[
\boxed{
\lim_{x\to\infty}
\left(
\sqrt{x^2+x}-x
\right)
=
\frac12.
}
\]
\end{answerbox}

\end{workedexamplebox}

%%%%%%%%%%%%%%%%%%%%%%%%%%%%%%%%%%%%%%%%%%%%%%%%%%%%%%%%%%%%
\subsection{Indeterminate Powers}
\label{subsec:indeterminate-powers-reference}
%%%%%%%%%%%%%%%%%%%%%%%%%%%%%%%%%%%%%%%%%%%%%%%%%%%%%%%%%%%%

The forms

\[
\boxed{
0^0,
\qquad
1^\infty,
\qquad
\infty^0
}
\]

are indeterminate.

A common method is logarithmic transformation.

If

\[
y=f(x)^{g(x)},
\]

then

\[
\boxed{
\ln y
=
g(x)\ln f(x).
}
\]

Analyze the logarithmic limit first, then exponentiate.

%%%%%%%%%%%%%%%%%%%%%%%%%%%%%%%%%%%%%%%%%%%%%%%%%%%%%%%%%%%%
\subsection{Worked Example: \(1^\infty\)}
\label{subsec:worked-one-infinity}
%%%%%%%%%%%%%%%%%%%%%%%%%%%%%%%%%%%%%%%%%%%%%%%%%%%%%%%%%%%%

\begin{workedexamplebox}{An Exponential Limit}

Evaluate

\[
\lim_{x\to\infty}
\left(
1+\frac{2}{x}
\right)^x.
\]

Use the standard limit

\[
\left(
1+\frac{a}{x}
\right)^x
\to
e^a.
\]

Therefore,

\begin{answerbox}
\[
\boxed{
\lim_{x\to\infty}
\left(
1+\frac{2}{x}
\right)^x
=
e^2.
}
\]
\end{answerbox}

\end{workedexamplebox}

%%%%%%%%%%%%%%%%%%%%%%%%%%%%%%%%%%%%%%%%%%%%%%%%%%%%%%%%%%%%
\subsection{Limit Comparison Through Ratios}
\label{subsec:ratio-comparison-limits}
%%%%%%%%%%%%%%%%%%%%%%%%%%%%%%%%%%%%%%%%%%%%%%%%%%%%%%%%%%%%

If

\[
\boxed{
\lim_{x\to a}
\frac{f(x)}{g(x)}
=
1,
}
\]

then

\[
\boxed{
f(x)\sim g(x).
}
\]

If instead

\[
\boxed{
\lim_{x\to a}
\frac{f(x)}{g(x)}
=
0,
}
\]

then \(f\) is asymptotically smaller than \(g\).

%%%%%%%%%%%%%%%%%%%%%%%%%%%%%%%%%%%%%%%%%%%%%%%%%%%%%%%%%%%%
\subsection{Dominant Balance}
\label{subsec:dominant-balance-reference}
%%%%%%%%%%%%%%%%%%%%%%%%%%%%%%%%%%%%%%%%%%%%%%%%%%%%%%%%%%%%

When several terms occur, the dominant terms often determine the limit.

For example,

\[
\frac{
7x^5-3x^2+1
}{
2x^5+x
}
\]

has leading terms

\[
7x^5
\]

and

\[
2x^5.
\]

Thus

\[
\boxed{
\frac{
7x^5-3x^2+1
}{
2x^5+x
}
\to
\frac72
}
\]

as \(|x|\to\infty\).

%%%%%%%%%%%%%%%%%%%%%%%%%%%%%%%%%%%%%%%%%%%%%%%%%%%%%%%%%%%%
\subsection{Limits and Inequalities}
\label{subsec:limits-and-inequalities}
%%%%%%%%%%%%%%%%%%%%%%%%%%%%%%%%%%%%%%%%%%%%%%%%%%%%%%%%%%%%

Suppose

\[
f(x)\leq g(x)
\]

near \(a\), and both limits exist.

Then

\[
\boxed{
\lim_{x\to a}f(x)
\leq
\lim_{x\to a}g(x).
}
\]

Order is preserved under limits.

%%%%%%%%%%%%%%%%%%%%%%%%%%%%%%%%%%%%%%%%%%%%%%%%%%%%%%%%%%%%
\subsection{Nonnegativity and Limits}
\label{subsec:nonnegative-limit-rule}
%%%%%%%%%%%%%%%%%%%%%%%%%%%%%%%%%%%%%%%%%%%%%%%%%%%%%%%%%%%%

If

\[
f(x)\geq0
\]

near \(a\) and

\[
f(x)\to L,
\]

then

\[
\boxed{
L\geq0.
}
\]

%%%%%%%%%%%%%%%%%%%%%%%%%%%%%%%%%%%%%%%%%%%%%%%%%%%%%%%%%%%%
\subsection{Uniqueness of Limits}
\label{subsec:uniqueness-of-limits}
%%%%%%%%%%%%%%%%%%%%%%%%%%%%%%%%%%%%%%%%%%%%%%%%%%%%%%%%%%%%

If a limit exists, it is unique.

A function cannot satisfy both

\[
\lim_{x\to a}f(x)=L
\]

and

\[
\lim_{x\to a}f(x)=M
\]

with

\[
L\neq M.
\]

%%%%%%%%%%%%%%%%%%%%%%%%%%%%%%%%%%%%%%%%%%%%%%%%%%%%%%%%%%%%
\subsection{Local Boundedness From a Finite Limit}
\label{subsec:finite-limit-local-boundedness}
%%%%%%%%%%%%%%%%%%%%%%%%%%%%%%%%%%%%%%%%%%%%%%%%%%%%%%%%%%%%

If

\[
\lim_{x\to a}f(x)=L
\]

with finite \(L\), then \(f\) is bounded in some punctured neighborhood of
\(a\).

Thus finite convergence prevents arbitrarily large local oscillation near
the point.

%%%%%%%%%%%%%%%%%%%%%%%%%%%%%%%%%%%%%%%%%%%%%%%%%%%%%%%%%%%%
\subsection{Limits Ignore the Function Value at the Point}
\label{subsec:limit-independent-point-value}
%%%%%%%%%%%%%%%%%%%%%%%%%%%%%%%%%%%%%%%%%%%%%%%%%%%%%%%%%%%%

The value

\[
f(a)
\]

does not directly determine

\[
\lim_{x\to a}f(x).
\]

For example, a function may satisfy

\[
f(a)=100
\]

while

\[
\lim_{x\to a}f(x)=3.
\]

Limits concern nearby behavior.

%%%%%%%%%%%%%%%%%%%%%%%%%%%%%%%%%%%%%%%%%%%%%%%%%%%%%%%%%%%%
\subsection{Removable Discontinuity}
\label{subsec:removable-discontinuity-limit}
%%%%%%%%%%%%%%%%%%%%%%%%%%%%%%%%%%%%%%%%%%%%%%%%%%%%%%%%%%%%

Suppose

\[
\lim_{x\to a}f(x)=L
\]

exists but either

\[
f(a)
\]

is undefined or

\[
f(a)\neq L.
\]

Then the discontinuity can be removed by defining

\[
\boxed{
f(a)=L.
}
\]

%%%%%%%%%%%%%%%%%%%%%%%%%%%%%%%%%%%%%%%%%%%%%%%%%%%%%%%%%%%%
\subsection{Jump Discontinuity}
\label{subsec:jump-discontinuity-limit}
%%%%%%%%%%%%%%%%%%%%%%%%%%%%%%%%%%%%%%%%%%%%%%%%%%%%%%%%%%%%

A jump occurs when finite one-sided limits exist but disagree:

\[
\boxed{
\lim_{x\to a^-}f(x)
\neq
\lim_{x\to a^+}f(x).
}
\]

Then the two-sided limit does not exist.

%%%%%%%%%%%%%%%%%%%%%%%%%%%%%%%%%%%%%%%%%%%%%%%%%%%%%%%%%%%%
\subsection{Infinite Discontinuity}
\label{subsec:infinite-discontinuity-limit}
%%%%%%%%%%%%%%%%%%%%%%%%%%%%%%%%%%%%%%%%%%%%%%%%%%%%%%%%%%%%

If

\[
f(x)
\to
\pm\infty
\]

from at least one side as \(x\to a\), the function has an infinite
discontinuity at \(a\).

A vertical asymptote often occurs.

%%%%%%%%%%%%%%%%%%%%%%%%%%%%%%%%%%%%%%%%%%%%%%%%%%%%%%%%%%%%
\subsection{Oscillatory Failure of a Limit}
\label{subsec:oscillatory-limit-failure}
%%%%%%%%%%%%%%%%%%%%%%%%%%%%%%%%%%%%%%%%%%%%%%%%%%%%%%%%%%%%

The function

\[
\boxed{
\sin\left(\frac1x\right)
}
\]

does not have a limit as

\[
x\to0.
\]

It oscillates infinitely often between \(-1\) and \(1\).

However,

\[
x\sin\left(\frac1x\right)
\to0
\]

by the squeeze theorem.

%%%%%%%%%%%%%%%%%%%%%%%%%%%%%%%%%%%%%%%%%%%%%%%%%%%%%%%%%%%%
\subsection{Common Special Limits}
\label{subsec:common-special-limits-table}
%%%%%%%%%%%%%%%%%%%%%%%%%%%%%%%%%%%%%%%%%%%%%%%%%%%%%%%%%%%%

\[
\boxed{
\lim_{x\to0}
\frac{\sin x}{x}
=
1
}
\]

\[
\boxed{
\lim_{x\to0}
\frac{\tan x}{x}
=
1
}
\]

\[
\boxed{
\lim_{x\to0}
\frac{1-\cos x}{x^2}
=
\frac12
}
\]

\[
\boxed{
\lim_{x\to0}
\frac{e^x-1}{x}
=
1
}
\]

\[
\boxed{
\lim_{x\to0}
\frac{\ln(1+x)}{x}
=
1
}
\]

\[
\boxed{
\lim_{x\to0}
(1+x)^{1/x}
=
e
}
\]

\[
\boxed{
\lim_{n\to\infty}
\left(
1+\frac1n
\right)^n
=
e
}
\]

%%%%%%%%%%%%%%%%%%%%%%%%%%%%%%%%%%%%%%%%%%%%%%%%%%%%%%%%%%%%
\subsection{Quick Limit-Law Table}
\label{subsec:quick-limit-law-table}
%%%%%%%%%%%%%%%%%%%%%%%%%%%%%%%%%%%%%%%%%%%%%%%%%%%%%%%%%%%%

If

\[
f(x)\to L,
\qquad
g(x)\to M,
\]

then:

\[
\boxed{
c
\to
c
}
\]

\[
\boxed{
cf(x)
\to
cL
}
\]

\[
\boxed{
f(x)+g(x)
\to
L+M
}
\]

\[
\boxed{
f(x)-g(x)
\to
L-M
}
\]

\[
\boxed{
f(x)g(x)
\to
LM
}
\]

\[
\boxed{
\frac{f(x)}{g(x)}
\to
\frac{L}{M},
\qquad
M\neq0
}
\]

\[
\boxed{
[f(x)]^n
\to
L^n
}
\]

\[
\boxed{
|f(x)|
\to
|L|.
}
\]

%%%%%%%%%%%%%%%%%%%%%%%%%%%%%%%%%%%%%%%%%%%%%%%%%%%%%%%%%%%%
\subsection{Common Errors}
\label{subsec:limit-laws-common-errors}
%%%%%%%%%%%%%%%%%%%%%%%%%%%%%%%%%%%%%%%%%%%%%%%%%%%%%%%%%%%%

\subsubsection{Treating \(0/0\) as Zero}

The expression

\[
\boxed{
\frac00
}
\]

is indeterminate.

It does not equal zero.

%%%%%%%%%%%%%%%%%%%%%%%%%%%%%%%%%%%%%%%%%%%%%%%%%%%%%%%%%%%%
\subsubsection{Treating \(\infty/\infty\) as One}

The form

\[
\boxed{
\frac{\infty}{\infty}
}
\]

does not imply a limit of \(1\).

For example,

\[
\frac{x}{x^2}
\to0,
\]

while

\[
\frac{x^2}{x}
\to\infty.
\]

%%%%%%%%%%%%%%%%%%%%%%%%%%%%%%%%%%%%%%%%%%%%%%%%%%%%%%%%%%%%
\subsubsection{Treating Infinity as an Ordinary Number}

Do not use ordinary arithmetic rules blindly with \(\infty\).

For example,

\[
\infty-\infty
\]

is indeterminate.

%%%%%%%%%%%%%%%%%%%%%%%%%%%%%%%%%%%%%%%%%%%%%%%%%%%%%%%%%%%%
\subsubsection{Assuming the Function Value Determines the Limit}

It is possible that

\[
f(a)
\neq
\lim_{x\to a}f(x).
\]

The function may even be undefined at \(a\) while the limit exists.

%%%%%%%%%%%%%%%%%%%%%%%%%%%%%%%%%%%%%%%%%%%%%%%%%%%%%%%%%%%%
\subsubsection{Ignoring One-Sided Limits}

A two-sided limit exists only if the left and right limits agree.

%%%%%%%%%%%%%%%%%%%%%%%%%%%%%%%%%%%%%%%%%%%%%%%%%%%%%%%%%%%%
\subsubsection{Canceling at the Point Instead of Near the Point}

When

\[
\frac{(x-a)(x+b)}{x-a}
\]

is simplified, cancellation is valid only for

\[
x\neq a.
\]

This is sufficient for the limit because limits use nearby values.

%%%%%%%%%%%%%%%%%%%%%%%%%%%%%%%%%%%%%%%%%%%%%%%%%%%%%%%%%%%%
\subsubsection{Forgetting \(\sqrt{x^2}=|x|\)}

This error is especially serious for

\[
x\to-\infty.
\]

For negative \(x\),

\[
|x|=-x.
\]

%%%%%%%%%%%%%%%%%%%%%%%%%%%%%%%%%%%%%%%%%%%%%%%%%%%%%%%%%%%%
\subsubsection{Using Trigonometric Fundamental Limits in Degrees}

The standard formulas

\[
\frac{\sin x}{x}\to1
\]

and

\[
\frac{1-\cos x}{x^2}\to\frac12
\]

require \(x\) in radians.

%%%%%%%%%%%%%%%%%%%%%%%%%%%%%%%%%%%%%%%%%%%%%%%%%%%%%%%%%%%%
\subsubsection{Using L'Hospital's Rule on a Nonindeterminate Form}

L'Hospital's Rule is designed for certain

\[
0/0
\]

and

\[
\infty/\infty
\]

forms.

It should not be applied automatically to every quotient.

%%%%%%%%%%%%%%%%%%%%%%%%%%%%%%%%%%%%%%%%%%%%%%%%%%%%%%%%%%%%
\subsubsection{Assuming a Fixed-Point Equation Proves Sequence Convergence}

From

\[
L=F(L)
\]

one obtains only candidate limits.

One must independently prove the sequence converges.

%%%%%%%%%%%%%%%%%%%%%%%%%%%%%%%%%%%%%%%%%%%%%%%%%%%%%%%%%%%%
\subsubsection{Confusing Limit at Infinity With Infinite Limit}

The statement

\[
\boxed{
\lim_{x\to\infty}f(x)=L
}
\]

describes input tending to infinity.

The statement

\[
\boxed{
\lim_{x\to a}f(x)=\infty
}
\]

describes output becoming unbounded.

These are different concepts.

%%%%%%%%%%%%%%%%%%%%%%%%%%%%%%%%%%%%%%%%%%%%%%%%%%%%%%%%%%%%
\subsubsection{Assuming Bounded Sequences Converge}

The sequence

\[
(-1)^n
\]

is bounded but divergent.

%%%%%%%%%%%%%%%%%%%%%%%%%%%%%%%%%%%%%%%%%%%%%%%%%%%%%%%%%%%%
\subsubsection{Assuming Oscillation Always Prevents a Product Limit}

Although

\[
\sin(1/x)
\]

does not converge as \(x\to0\),

\[
x\sin(1/x)
\to0.
\]

A shrinking factor can dominate bounded oscillation.

%%%%%%%%%%%%%%%%%%%%%%%%%%%%%%%%%%%%%%%%%%%%%%%%%%%%%%%%%%%%
\subsection{Limit Evaluation Workflow}
\label{subsec:limit-evaluation-workflow}
%%%%%%%%%%%%%%%%%%%%%%%%%%%%%%%%%%%%%%%%%%%%%%%%%%%%%%%%%%%%

A practical order of operations is:

\begin{enumerate}
    \item identify the limiting point or direction;
    \item try direct substitution;
    \item if a finite value results, stop when continuity justifies it;
    \item if \(0/0\) occurs, factor or rationalize;
    \item use trigonometric identities when needed;
    \item use fundamental limits;
    \item consider the squeeze theorem for bounded oscillation;
    \item for limits at infinity, identify dominant terms;
    \item divide rational functions by the largest relevant power;
    \item use \(\sqrt{x^2}=|x|\) when radicals are involved;
    \item transform other indeterminate forms algebraically;
    \item use L'Hospital's Rule only when applicable;
    \item use series or asymptotic equivalence for advanced limits;
    \item check one-sided behavior when signs or discontinuities matter.
\end{enumerate}

%%%%%%%%%%%%%%%%%%%%%%%%%%%%%%%%%%%%%%%%%%%%%%%%%%%%%%%%%%%%
\subsection{Exercises}
\label{subsec:limit-laws-exercises}
%%%%%%%%%%%%%%%%%%%%%%%%%%%%%%%%%%%%%%%%%%%%%%%%%%%%%%%%%%%%

\begin{exercise}[1]
Evaluate

\[
\lim_{x\to3}(2x+5).
\]
\end{exercise}

\begin{exercise}[1]
Evaluate

\[
\lim_{x\to-2}(x^2+4x+1).
\]
\end{exercise}

\begin{exercise}[1]
Evaluate

\[
\lim_{x\to1}
\frac{x+2}{x+5}.
\]
\end{exercise}

\begin{exercise}[1]
State the quotient limit law.
\end{exercise}

\begin{exercise}[1]
State the condition required on the denominator limit in the quotient law.
\end{exercise}

\begin{exercise}[1]
Explain what

\[
\lim_{x\to a^-}f(x)
\]

means.
\end{exercise}

\begin{exercise}[1]
Explain what

\[
\lim_{x\to a^+}f(x)
\]

means.
\end{exercise}

\begin{exercise}[2]
Evaluate

\[
\lim_{x\to4}
\frac{x^2-16}{x-4}.
\]
\end{exercise}

\begin{exercise}[2]
Evaluate

\[
\lim_{x\to-3}
\frac{x^2+5x+6}{x+3}.
\]
\end{exercise}

\begin{exercise}[2]
Evaluate

\[
\lim_{x\to9}
\frac{\sqrt{x}-3}{x-9}.
\]
\end{exercise}

\begin{exercise}[2]
Evaluate

\[
\lim_{x\to0}
\frac{\sin7x}{x}.
\]
\end{exercise}

\begin{exercise}[2]
Evaluate

\[
\lim_{x\to0}
\frac{\tan4x}{3x}.
\]
\end{exercise}

\begin{exercise}[2]
Evaluate

\[
\lim_{x\to0}
\frac{1-\cos5x}{x^2}.
\]
\end{exercise}

\begin{exercise}[2]
Evaluate

\[
\lim_{x\to0}
\frac{e^{3x}-1}{x}.
\]
\end{exercise}

\begin{exercise}[2]
Evaluate

\[
\lim_{x\to0}
\frac{\ln(1+4x)}{x}.
\]
\end{exercise}

\begin{exercise}[2]
Evaluate using the squeeze theorem:

\[
\lim_{x\to0}
x^3\cos\left(\frac1x\right).
\]
\end{exercise}

\begin{exercise}[2]
Determine whether

\[
\lim_{x\to0}
\sin\left(\frac1x\right)
\]

exists.
\end{exercise}

\begin{exercise}[2]
Evaluate the left- and right-hand limits of

\[
\frac1x
\]

at \(x=0\).
\end{exercise}

\begin{exercise}[2]
Evaluate

\[
\lim_{x\to\infty}
\frac{3x^2+1}{5x^2-7}.
\]
\end{exercise}

\begin{exercise}[2]
Evaluate

\[
\lim_{x\to\infty}
\frac{2x+1}{x^3+4}.
\]
\end{exercise}

\begin{exercise}[2]
Determine the end behavior of

\[
\frac{x^4+1}{x^2+3}.
\]
\end{exercise}

\begin{exercise}[2]
Evaluate

\[
\lim_{x\to-\infty}
\frac{\sqrt{x^2+1}}{x}.
\]
\end{exercise}

\begin{exercise}[2]
Find the horizontal asymptote of

\[
f(x)
=
\frac{4x^2+1}{2x^2-3}.
\]
\end{exercise}

\begin{exercise}[2]
Determine the vertical asymptote behavior of

\[
f(x)
=
\frac{1}{x-2}.
\]
\end{exercise}

\begin{exercise}[2]
Evaluate

\[
\lim_{n\to\infty}
\frac1{n^3}.
\]
\end{exercise}

\begin{exercise}[2]
Determine whether

\[
\left(
\frac23
\right)^n
\]

converges.
\end{exercise}

\begin{exercise}[2]
Determine whether

\[
(-1)^n
\]

converges.
\end{exercise}

\begin{exercise}[3]
Use subsequences to prove that

\[
(-1)^n
\]

does not converge.
\end{exercise}

\begin{exercise}[3]
Suppose

\[
a_{n+1}
=
\sqrt{3+a_n}.
\]

Assuming convergence, find all possible limits.
\end{exercise}

\begin{exercise}[3]
Explain why solving

\[
L=F(L)
\]

does not prove a recursive sequence converges.
\end{exercise}

\begin{exercise}[3]
Evaluate

\[
\lim_{x\to0}
\frac{\sin3x}{\sin5x}.
\]
\end{exercise}

\begin{exercise}[3]
Evaluate

\[
\lim_{x\to0}
\frac{\tan2x}{\sin7x}.
\]
\end{exercise}

\begin{exercise}[3]
Evaluate

\[
\lim_{x\to0}
\frac{e^{2x}-1}{\ln(1+3x)}.
\]
\end{exercise}

\begin{exercise}[3]
Evaluate

\[
\lim_{x\to0}
\frac{
1-\cos x
}{
\sin^2x
}.
\]
\end{exercise}

\begin{exercise}[3]
Evaluate

\[
\lim_{x\to\infty}
\left(
\sqrt{x^2+3x}-x
\right).
\]
\end{exercise}

\begin{exercise}[3]
Evaluate

\[
\lim_{x\to\infty}
\frac{\ln x}{x}.
\]
\end{exercise}

\begin{exercise}[3]
Evaluate

\[
\lim_{x\to\infty}
\frac{x^5}{e^x}.
\]
\end{exercise}

\begin{exercise}[3]
Evaluate

\[
\lim_{n\to\infty}
\frac{3^n}{n!}.
\]
\end{exercise}

\begin{exercise}[3]
Evaluate

\[
\lim_{n\to\infty}
\left(
1+\frac4n
\right)^n.
\]
\end{exercise}

\begin{exercise}[3]
Evaluate

\[
\lim_{x\to0}
(1+3x)^{1/x}.
\]
\end{exercise}

\begin{exercise}[3]
Use L'Hospital's Rule to evaluate

\[
\lim_{x\to0}
\frac{\sin x}{x}.
\]
\end{exercise}

\begin{exercise}[3]
Use L'Hospital's Rule to evaluate

\[
\lim_{x\to\infty}
\frac{\ln x}{x}.
\]
\end{exercise}

\begin{exercise}[3]
Use L'Hospital's Rule to evaluate

\[
\lim_{x\to0}
\frac{1-\cos x}{x^2}.
\]
\end{exercise}

\begin{exercise}[3]
Convert a \(0\cdot\infty\) expression into a quotient and evaluate it.
\end{exercise}

\begin{exercise}[3]
Convert an \(\infty-\infty\) expression into a form suitable for evaluation.
\end{exercise}

\begin{exercise}[3]
Use logarithms to evaluate an indeterminate power of type \(1^\infty\).
\end{exercise}

\begin{exercise}[3]
Show that

\[
\sin x\sim x
\]

as \(x\to0\).
\end{exercise}

\begin{exercise}[3]
Show that

\[
e^x-1\sim x
\]

as \(x\to0\).
\end{exercise}

\begin{exercise}[3]
Show that

\[
\ln(1+x)\sim x
\]

as \(x\to0\).
\end{exercise}

\begin{exercise}[3]
Show that

\[
1-\cos x
\sim
\frac{x^2}{2}
\]

as \(x\to0\).
\end{exercise}

\begin{exercise}[3]
Use a Taylor expansion to evaluate

\[
\lim_{x\to0}
\frac{
\sin x-x
}{
x^3
}.
\]
\end{exercise}

\begin{exercise}[3]
Use a Taylor expansion to evaluate

\[
\lim_{x\to0}
\frac{
\ln(1+x)-x
}{
x^2
}.
\]
\end{exercise}

\begin{exercise}[4]
Prove the uniqueness of limits using the epsilon--delta definition.
\end{exercise}

\begin{exercise}[4]
Prove the sum law for limits from the epsilon--delta definition.
\end{exercise}

\begin{exercise}[4]
Prove the product law for limits.
\end{exercise}

\begin{exercise}[4]
Prove the quotient law under the assumption that the denominator limit is
nonzero.
\end{exercise}

\begin{exercise}[4]
Prove the squeeze theorem.
\end{exercise}

\begin{exercise}[4]
Give an epsilon--delta proof that

\[
\lim_{x\to2}
(3x+1)
=
7.
\]
\end{exercise}

\begin{exercise}[4]
Give an epsilon--delta proof that

\[
\lim_{x\to1}
x^2
=
1.
\]
\end{exercise}

\begin{exercise}[4]
Prove that finite limits imply local boundedness.
\end{exercise}

\begin{exercise}[4]
Prove that limits preserve inequalities.
\end{exercise}

\begin{exercise}[4]
Prove that every convergent sequence is bounded.
\end{exercise}

\begin{exercise}[4]
Prove the monotone convergence theorem for real sequences using completeness.
\end{exercise}

\begin{exercise}[4]
Study the Bolzano--Weierstrass theorem for bounded sequences.
\end{exercise}

\begin{exercise}[4]
Derive

\[
\lim_{x\to0}
\frac{\sin x}{x}=1
\]

geometrically.
\end{exercise}

\begin{exercise}[4]
Derive

\[
\lim_{x\to0}
\frac{1-\cos x}{x^2}
=
\frac12
\]

from the sine limit.
\end{exercise}

\begin{exercise}[4]
Prove

\[
\lim_{x\to\infty}
\frac{x^n}{e^x}
=
0
\]

for positive integer \(n\).
\end{exercise}

\begin{exercise}[4]
Compare logarithmic, polynomial, exponential, and factorial growth using
limits.
\end{exercise}

\begin{exercise}[4]
Study asymptotic equivalence as an equivalence relation on suitable positive
functions.
\end{exercise}

%%%%%%%%%%%%%%%%%%%%%%%%%%%%%%%%%%%%%%%%%%%%%%%%%%%%%%%%%%%%
\subsection{Conceptual Summary}
\label{subsec:limit-laws-summary}
%%%%%%%%%%%%%%%%%%%%%%%%%%%%%%%%%%%%%%%%%%%%%%%%%%%%%%%%%%%%

The central limit laws allow limits to pass through algebraic operations.

If

\[
f(x)\to L
\]

and

\[
g(x)\to M,
\]

then

\[
\boxed{
f+g\to L+M,
}
\]

\[
\boxed{
fg\to LM,
}
\]

and

\[
\boxed{
\frac{f}{g}
\to
\frac{L}{M}
}
\]

when

\[
M\neq0.
\]

Direct substitution works for continuous functions.

When substitution produces

\[
\boxed{
\frac00,
}
\]

one should simplify before drawing conclusions.

Fundamental local limits include

\[
\boxed{
\frac{\sin x}{x}\to1,
}
\]

\[
\boxed{
\frac{1-\cos x}{x^2}\to\frac12,
}
\]

\[
\boxed{
\frac{e^x-1}{x}\to1,
}
\]

and

\[
\boxed{
\frac{\ln(1+x)}{x}\to1.
}
\]

Limits at infinity are often controlled by dominant growth terms.

For rational functions, compare polynomial degrees.

For sequences, convergence means

\[
\boxed{
a_n\to L
}
\]

as \(n\to\infty\).

Asymptotic notation such as

\[
\boxed{
f\sim g,
\qquad
f=O(g),
\qquad
f=o(g)
}
\]

describes relative growth.

The most important practical rules are:

\[
\boxed{
\text{do not treat indeterminate forms as answers},
}
\]

\[
\boxed{
\text{do not treat infinity as an ordinary real number},
}
\]

and

\[
\boxed{
\text{always inspect one-sided behavior when discontinuities are possible}.
}
\]

%%%%%%%%%%%%%%%%%%%%%%%%%%%%%%%%%%%%%%%%%%%%%%%%%%%%%%%%%%%%
\subsection{Section Connections}
\label{subsec:limit-laws-connections}
%%%%%%%%%%%%%%%%%%%%%%%%%%%%%%%%%%%%%%%%%%%%%%%%%%%%%%%%%%%%

\chapterconnection{
Limit Laws provides the reference foundation for continuity, differentiation,
integration, infinite series, asymptotic analysis, probability limits,
dynamical systems, and numerical approximation. Algebraic simplification
connects limits with Section~\ref{sec:algebraic-identities}; fundamental
trigonometric limits connect with Section~\ref{sec:trigonometric-identities};
sequence convergence connects with Analysis and Probability; limits at
infinity connect with asymptotics; and continuity allows direct substitution
through many standard functions. The next reference section, Derivative
Tables, collects derivative formulas for algebraic, exponential, logarithmic,
trigonometric, inverse-trigonometric, hyperbolic, inverse-hyperbolic,
multivariable, and commonly composed functions.
}

%%%%%%%%%%%%%%%%%%%%%%%%%%%%%%%%%%%%%%%%%%%%%%%%%%%%%%%%%%%%
% SECTION SOURCE TRACKING
%%%%%%%%%%%%%%%%%%%%%%%%%%%%%%%%%%%%%%%%%%%%%%%%%%%%%%%%%%%%

%------------------------------------------------------------
% 14.8 Limit Laws
%------------------------------------------------------------
%
% Source basis:
%   - Standard single-variable calculus.
%   - Standard real analysis and sequence convergence.
%   - Standard asymptotic and elementary limit techniques.
%
% Used for:
%   - finite limits;
%   - epsilon-delta definition;
%   - one-sided limits;
%   - two-sided limits;
%   - constant law;
%   - identity law;
%   - constant-multiple law;
%   - sum and difference laws;
%   - product and quotient laws;
%   - power and root laws;
%   - absolute-value law;
%   - polynomial limits;
%   - rational-function limits;
%   - direct substitution;
%   - indeterminate forms;
%   - factoring;
%   - rationalization;
%   - piecewise limits;
%   - squeeze theorem;
%   - fundamental trigonometric limits;
%   - exponential and logarithmic limits;
%   - composition;
%   - continuity;
%   - limits at infinity;
%   - horizontal asymptotes;
%   - vertical asymptotes;
%   - dominant terms;
%   - rational functions at infinity;
%   - radical limits;
%   - infinite limits;
%   - growth hierarchy;
%   - sequence limits;
%   - geometric sequences;
%   - monotone convergence;
%   - subsequences;
%   - recursive sequences;
%   - asymptotic equivalence;
%   - Big-O notation;
%   - little-o notation;
%   - Taylor-based limit evaluation;
%   - L'Hospital's Rule;
%   - transformed indeterminate forms;
%   - limit comparison;
%   - uniqueness;
%   - discontinuity types;
%   - common errors;
%   - applications and exercises.
%
% Notes:
%   - Derivative Tables is developed next in Section 14.9.
%   - Integral Tables follows in Section 14.10.
%
%%%%%%%%%%%%%%%%%%%%%%%%%%%%%%%%%%%%%%%%%%%%%%%%%%%%%%%%%%%%